% Generated from locale/tr/content/first-order-logic/natural-deduction/rules-and-proofs.tex; do not hand-edit.


\iftag{FOL}
      {\olfileid[tr]{fol}{ntd}{rul}}
      {\olfileid[tr]{pl}{ntd}{rul}}

\olsection{Kurallar ve \usetoken{P}{derivation}}

\begin{explain}
Doğal tümdengelim sistemleri, matematiksel ispatlarda kullanılan biçimsel
olmayan akıl yürütmeye yakından koşut olacak şekilde tasarlanmıştır (bu
nedenle bir ölçüde ``doğal''dır). Doğal tümdengelim ispatları varsayımlarla
başlar. Ardından çıkarım kuralları uygulanır. Varsayımlar, \Intro{\lnot},
\Intro{\lif}, \iftag{FOL}{\Elim{\lor} ve \Elim{\lexists}}{ ve
\Elim{\lor}} çıkarım kurallarıyla ``!!{discharged}'' sayılır; açıklık sağlamak için
!!{discharged} varsayımın etiketi çıkarımın yanına yazılır.
\end{explain}

\begin{defn}[Varsayım]
Bir \emph{varsayım}, herhangi bir dalın en üst konumundaki herhangi bir
!!{sentence} olarak alınır.
\end{defn}

Doğal tümdengelimde !!^{derivation}s, !!{sentence}s içeren belirli ağaçlardır;
en üstteki !!{sentence}s varsayımlardır ve !!a{sentence}, bir, iki veya üç
başka tümcenin altında duruyorsa bir çıkarım kuralına doğru biçimde dayanmalıdır.
Çıkarımın üstündeki !!{sentence}s \emph{öncüller}, altındaki !!{sentence} ise
çıkarımın \emph{sonucu} diye adlandırılır. Kurallar çiftler hâlinde gelir: her
!!{operator} için bir giriş ve bir giderme kuralı vardır. Bunlar sonuçta
!!a{operator} ortaya çıkarır veya kuralın bir öncülünden !!a{operator}
kaldırır. Bazı kurallar, belirli türden bir varsayımın
\emph{!!{discharged}} olmasına izin verir. Hangi varsayımın hangi çıkarımla
!!{discharged} olduğunu belirtmek için hem varsayıma hem de çıkarıma etiket
veririz. Bu, varsayımın ``$\Discharge{!A}{n}$'' biçiminde yazılmasıyla
gösterilir.

% Only include this sentence is on of \land, lor, lif or lnot is defined.

\iftag{notprvNot,notprvAnd,notprvOr,notprvIf}{}%
	{Bazıları tanımlı olsa bile bütün !!{operator}s, yani $\land$, $\lor$,
  $\lif$, $\lnot$ ve $\lfalse$ için kuralların ele alınması gelenekseldir.}



